\section{Bối cảnh nghiên cứu và Thực trạng ngành}

Trong kỷ nguyên chuyển đổi số và bùng nổ của thương mại điện tử, ngành du lịch và vận tải hành khách tại Việt Nam đang chứng kiến sự dịch chuyển mạnh mẽ từ mô hình giao dịch trực tiếp sang các nền tảng trực tuyến đa kênh (Omnichannel). Người dùng hiện đại, đặc biệt là thế hệ trẻ, có xu hướng tự lên kế hoạch cho các chuyến đi tự túc (Free \& Easy) và kỳ vọng nhận được sự hỗ trợ tư vấn tức thời thông qua các ứng dụng nhắn tin tức thời như Facebook Messenger, Zalo OA và các tiện ích Webchat.

Đặc biệt, nghiệp vụ du lịch liên tỉnh tại các thị trường trọng điểm (ví dụ: các tuyến vận tải hành khách kết hợp lưu trú như Sài Gòn -- Đà Lạt, Sài Gòn -- Nha Trang) có tính liên kết chặt chẽ giữa hai dịch vụ cốt lõi:
\begin{enumerate}
    \item \textbf{Dịch vụ vận tải hành khách:} Chủ yếu là các dòng xe khách giường nằm chất lượng cao (Limousine 34 phòng, Cabin đôi) với lịch trình xuất phát tập trung vào ban đêm (từ 21:00 đến 23:30) nhằm giúp hành khách tiết kiệm thời gian di chuyển.
    \item \textbf{Dịch vụ lưu trú du lịch:} Bao gồm hệ thống khách sạn, resort và homestay với các quy định khắt khe về thời gian nhận phòng tiêu chuẩn (thường là 14:00) và trả phòng (12:00 hôm sau), đi kèm các chính sách phụ thu nhận phòng sớm (Early Check-in), chính sách hoàn hủy (Cancellation Policy) và quy định lưu trú đặc thù (cho phép mang thú cưng, phụ thu giường phụ cho trẻ em).
\end{enumerate}

Sự bùng nổ của các Mô hình Ngôn ngữ Lớn (Large Language Models -- LLMs) như GPT-4, Gemini, Claude và các kiến trúc Tác tử Tự trị (Autonomous Agents) đã mở ra tiềm năng to lớn trong việc tự động hóa hoàn toàn quy trình hội thoại hướng tác vụ (Task-Oriented Dialogue Systems). Doanh nghiệp kỳ vọng có thể thay thế đội ngũ nhân viên trực tổng đài ca đêm bằng các trợ lý ảo AI có khả năng giao tiếp tự nhiên, thấu hiểu ngữ cảnh và chốt đơn tự động 24/7.
